\documentclass{article}

\usepackage{amsmath}
\usepackage{amssymb}
\usepackage{comment}


\usepackage[table]{xcolor}
\usepackage{longtable}
\usepackage{geometry}
\geometry{margin=1in}
\setcounter{secnumdepth}{0}
\setlength{\parindent}{0pt}
\usepackage{hyperref}
\usepackage{xcolor}
\hypersetup{colorlinks, linkcolor={red!50!black},
    citecolor={blue!50!black},
    urlcolor={blue!80!black}}

% Logo on header
\usepackage{graphicx}
\usepackage{fancyhdr}
\renewcommand{\headrulewidth}{1pt}
\fancypagestyle{plain}{%
  \fancyhead[L]{
    \begin{tabular}{ll}
      \begin{tabular}[t]{c}
        \includegraphics[scale=0.06]{UdeM.png}%
      \end{tabular} &
      \begin{tabular}[b]{l}
        Faculté des arts et des sciences\tabularnewline
        Département d'\textbf{informatique et de recherche opérationnelle}  \tabularnewline
      \end{tabular}
    \end{tabular}
  }%
}

% Title
\title{\textbf{Pretrained Language Models on Low Ressources}}
\author{Gao Yinghan, Wei Wei, Frederic Boileau}

\begin{document}
\maketitle
\thispagestyle{plain}


\section*{Introduction to attention and its history}

It was conjectured by Bahdanau et al\cite{bahdanau2016neural} that the
dominant encoder decoder paradigms to Neural Machine Translation presented one
major drawback. All the information necessary to the transduction mechanism had
to be encoded into a fixed length vector. The dominant approach was to use a RNN
to encode the sequence to transduce (or more precisely the



sentence to translate
in the case of Bahdanau et al) into a fixed length vector c which was a function
of the encoder's hiddens states, usually simply extracting the last hidden
parameter. To remedy the situation the authors tackled the problem with an
innovative approach: making the context vector depend on a linear combinations
of \textbf{all} the hidden states, where the weights assigned respectively were
learned through training a feedforward network. This enabled to effectively
learn what part of the input sentence the output had to \textit{attend to}. This
idea was called attention.

\begin{equation}
  \label{eq:3}
  c_{i} = \sum_{j=1}^{T_{x}}\alpha_{ij}h_{j} \qquad
  \alpha_{ij} = \frac{\exp(e_{ij})}{\sum_{k=1}^{T_{x}}\exp(e_{ik})} \qquad
  e_{ij} = a(s_{i-1}, h_{j})
\end{equation}
% TODO explain in words fundamental equations of attention


\section*{The Transformer revolution}
% TODO explain transformer
Later on Vaswani et al \cite{allyouneed} introduced a revolutionnary
architecture, the transformer. Still based on the idea of an encoder and decoder
for transduction, they based their architecture solely on an attention mechanism
enables to model the long range dependencies in a sentence without the
inherently sequential constraint which RNNs in their various forms imply. The
massive parallelization enabled massive improvements in training time to achieve
state of the art results in machine translation\\

% TODO add fundamental transformer equation
%

\section*{Pretrained Language Models}
% TODO explain PLMs in general and knowledge transfer
We will focus in our project not precisely on transduction tasks but on large
pretrained language models (PLMs) which can then be fined tuned for downstream
tasks such as \ldots\\



% TODO explain BERT
%
BERT is \ldots\cite{bert}\\


\begin{table}[h]
  \centering
  \begin{tabular}{cc}
    Architecture & Number of parameters\\
    \hline
    BERT & 340M\\
    GPT-2 & 1.5B\\
    MegatronLM & 8.3B\\
    T5 & 11B\\
    T-NLG & 17B\\
    GShard & 600B
  \end{tabular}
  \caption{PLMs and their sizes}
\end{table}

\section*{Compression approaches}
Model compression consists of \ldots

\subsection*{Knowledge distillation}
% TODO explain Knowledge distillation
%
Knowledge Distillation (KD)\cite{hinton2015distilling}

\subsection*{tinyBERT}
% TODO explain TinyBERT
%
TinyBERT\cite{tinybert}

\subsection*{Other model commpression techniques}
% TODO explain other model compression techniques
%
We have the following other model compression techniques\cite{gupta2020compression}
\begin{itemize}
  \item Pruning
  \item Quantization
  \item Parameter Sharing
  \item Tensor Decomposition
\end{itemize}



% TODO explain metrics such for evaluation such as Squad and Glue
%
PLMs can be evaluated on GLUE for general LM performance and more specific downstream
tasks for which they are fine tuned such as Squad.
\begin{itemize}
  \item SquAD\cite{squad}
  \item GLUE\cite{glue}
\end{itemize}


\medskip

\bibliography{projectbibliography}
\bibliographystyle{unsrt}



\begin{comment}

\begin{itemize}
  \item Albert\cite{Lan2020ALBERT}
  \item Reformer\cite{Kitaev2020Reformer:}
  \item Simple Recurrent Units (SRU)\cite{sru}
  \item Sparse Transformers (ST)\cite{sparsetransformers}
\end{itemize}

\end{comment}

\end{document}
